\documentclass[11pt]{ctexart}
\usepackage[margin=1in]{geometry}
\begin{document}
\section{导数的定义}
设函数 $f(x)$ 在点 $x_0$ 的某个邻域内有定义。若极限
\[ f'(x_0)=\lim_{h\to 0}\frac{f(x_0+h)-f(x_0)}{h} \]
存在，则称 $f$ 在 $x_0$ 处可导，该极限值称为导数。

\subsection{可导与连续}
可导必连续，连续不一定可导。例如 $y=|x|$ 在 $x=0$ 处连续但不可导。

\newpage
\section{链式法则}
复合函数的求导法则为
\[ (f\circ g)'(x)=f'(g(x))\cdot g'(x). \]
本页用于验证第二页能被单独提取，且页码是真实分页结果。
\end{document}
